% Network-Attached FPGA Inference Scanning: A Bit-Exact SAN Catastrophe Scan
% over 100G Ethernet, and a Cautionary Tale About ARP
%
% TechRxiv preprint submission
%
% Compile with: tectonic paper_c_san_net_fpga.tex
%           or: pdflatex paper_c_san_net_fpga.tex (run twice for refs/TOC)

\documentclass[11pt,a4paper]{article}

\usepackage[utf8]{inputenc}
\usepackage[T1]{fontenc}
\usepackage{lmodern}
\usepackage[margin=1in]{geometry}
\usepackage{amsmath,amssymb}
\usepackage{booktabs}
\usepackage{hyperref}
\usepackage{url}
\usepackage{xcolor}
\usepackage{listings}
\usepackage{caption}
\usepackage{microtype}
\usepackage{enumitem}

\hypersetup{
  colorlinks=true,
  linkcolor=blue!50!black,
  citecolor=blue!50!black,
  urlcolor=blue!50!black,
  pdftitle={Network-Attached FPGA Inference Scanning},
  pdfauthor={Demetrios Chiuratto Agourakis}
}

\lstset{
    basicstyle=\ttfamily\small,
    breaklines=true,
    frame=single,
    backgroundcolor=\color{gray!10},
    columns=fullflexible,
    keepspaces=true
}

\newcommand{\code}[1]{\texttt{#1}}

\title{\Large Network-Attached FPGA Inference Scanning: \\[4pt]
  A Bit-Exact SAN Catastrophe Scan over 100G Ethernet, \\[4pt]
  and a Cautionary Tale About ARP}

\author{
  Demetrios Chiuratto Agourakis\\
  \small ORCID: 0009-0001-8671-8878
}
\date{2026-08-13}

\begin{document}
\maketitle

\begin{abstract}
We port a previously-accepted, production FPGA kernel --- a
catastrophe-exit scan for early-exit inference architectures
(SAN-ResNet-50, SAN-ViT-large), validated bit-exact over PCIe DMA at
511~Msamples/s on an AMD Alveo U250 (spec T3, verdict \code{T3\_GREEN})
--- to receive its input directly from a 100G Ethernet fibre instead of
host DMA. The kernel's arithmetic is unchanged; only the data-delivery
interface changes, from an AXI4 memory-mapped read to an AXI4-Stream fed
by a third-party UDP/IP network stack (Xilinx/AMD's VNx). We report two
results. First, a correctness result: the network-attached kernel
reproduces the DMA-attached kernel's output bit-for-bit --- histogram,
catastrophe count, and accumulated MAC cost --- across cohorts up to
4{,}000{,}003 samples, confirmed in a batch of ten consecutive,
independent hardware runs with zero data loss. Second, a
systems-debugging result: an initial throughput measurement attempt
produced what appeared to be hardware instability --- a third-party
network IP core that would intermittently stop accepting traffic
regardless of the offered send rate, recoverable only by a full bitstream
reload. We show this was not a hardware fault at all: it was ARP cache
staleness in the host's Linux IP stack, which silently discards or delays
the first datagram of a burst sent after any idle period. A three-packet
ICMP warm-up immediately before each data burst eliminates the failure
entirely (10/10 clean runs post-fix, including a deliberately adversarial
run with the ARP entry forcibly deleted). We report this debugging
narrative in full because the failure signature --- a third-party
accelerator that ``just needs a reset'' under unpredictable conditions
--- is a common and expensive misdiagnosis in FPGA/network co-design, and
the actual root cause was one layer up the stack from where we first
looked.
\end{abstract}

\noindent\textbf{Keywords:} FPGA, Alveo U250, 100G Ethernet, RoCE fabric,
network-attached accelerators, early-exit inference, catastrophe scan,
VNx, HLS, bit-exact validation, ARP, systems debugging

\bigskip
\noindent\textbf{Category:} Computer Science $>$ Distributed, Parallel,
and Cluster Computing / Hardware Architecture

\noindent\textbf{License:} CC BY 4.0

\noindent\textbf{Status:} preprint draft, 2026-08-13. Companion artifact
repository: \code{fpga/san-net/} (kernel, host controller, injector, test
harness, and a complete build/debug log preserved in \code{README.md}).

\section{Introduction}

Early-exit inference architectures (SAN-ResNet-50, SAN-ViT-large) commit
to a prediction at the first internal stage whose confidence clears a
threshold, avoiding the compute cost of deeper stages. The
\emph{catastrophe scan} kernel implements this decision as a priority
encoder over quantized Q0.15 confidences, with MAC-cost metering that
charges each sample the real, architecture-specific FLOP prefix of its
exit point --- a host-programmable lookup table.

This kernel was previously accepted onto an AMD Alveo U250 (spec T3), fed
by host DMA over PCIe, and validated bit-exact against a software golden
model at 511~Msamples/s. That result establishes the kernel's
correctness and its DMA-bound throughput ceiling. It leaves open a
different question: what if the confidences do not originate on the host
at all, but arrive over the network from a producer elsewhere in a
cluster --- a GPU inference node, for instance, streaming logits to an
FPGA scanner as a disaggregated pipeline stage? This motivates a
\textbf{network-attached} variant: the same kernel, unchanged, fed by a
100G Ethernet link instead of DMA.

We build this variant on top of VNx (\code{xup\_vitis\_network\_example}),
AMD's reference UDP/IP network stack for the U250, and report what we
found in two parts: the correctness result (Section~\ref{sec:correctness}),
and the throughput measurement attempt that initially looked like a
hardware reliability problem and turned out to be a one-line host
networking fix (Section~\ref{sec:throughput}). We include the latter at
the same level of detail as the former because, in our reading of the
FPGA-networking literature and vendor documentation, this specific
failure mode --- third-party network IP appearing to wedge under load,
with the true cause sitting in the host's neighbor-discovery cache
rather than the fabric --- is under-documented, and we believe the
debugging method is as reusable as the result.

\section{Background}

\paragraph{The accepted kernel (spec T3).} \code{krnl\_san\_scan}
implements: (i) a catastrophe scan --- for each sample, a comparator
tree finds the first exit point whose confidence clears a threshold
\code{q\_delta}; samples that clear no threshold propagate to the final
head and are counted as catastrophes; and (ii) FLOP metering --- each
sample is charged the exact MAC-prefix cost of its exit point from a
host-loaded 64-bit lookup table, accumulated as a 64-bit total.
Confidences arrive pre-quantized (Q0.15, no floating point on the card).
The DMA-attached variant packs four samples per 512-bit AXI beat
(\code{LANES=4}) and is throughput-bound by the PCIe/DMA path at
511~Msamples/s, not by kernel logic (\code{II=1}).

\paragraph{The Alveo U250 and VNx.} The U250 is a Xilinx/AMD Virtex
UltraScale+ VU13P card (4 SLRs, 1.72M LUTs, 12{,}288 DSPs) with dual
QSFP28 cages, each capable of 100G Ethernet via a hardened CMAC block.
VNx is AMD's reference implementation of a UDP/IP network stack for
this platform: a \code{cmac} kernel (the Ethernet MAC/PCS, including
RS-FEC), a \code{networklayer} kernel (ARP responder, ICMP responder,
and a fixed-size UDP socket table that routes matching
\code{(theirIP, theirPort, myPort)} triples to per-socket AXI-Stream
ports), and helper DMA kernels (\code{krnl\_mm2s}/\code{krnl\_s2mm}) for
moving stream data to/from card DRAM. The \code{networklayer}'s UDP
socket table supports up to \code{MAX\_SOCKETS} entries (16 in our
build), each entry occupying an 8-byte stride even though its
constituent fields are 4 bytes --- a detail that, if assumed rather than
read from the generated kernel descriptor, silently misconfigures socket
routing.

\section{Architecture}

\subsection{Kernel changes}

The kernel's arithmetic core is untouched. Four changes port it to a
streaming interface:

\begin{enumerate}[label=(\arabic*)]
  \item The input argument becomes
    \code{hls::stream<ap\_axiu<512,0,0,16>>} in place of a
    \code{bus\_t*} pointer with an explicit index.
  \item The per-beat read becomes \code{samples\_in.read().data} in
    place of \code{samples[w]}.
  \item \code{\#pragma HLS INTERFACE axis port=samples\_in} replaces the
    \code{m\_axi} interface pragma for that argument.
  \item The 16-bit \code{TDEST} field is retained to match VNx's
    application-side stream convention, which uses \code{TDEST} to
    demultiplex per-socket traffic.
\end{enumerate}

Sample packing is unchanged and remains the compatibility contract with
the accepted artifact: one sample is 128 bits, holding up to seven
15-bit Q0.15 confidence fields, field $k$ occupying bits
$[15k{+}14 : 15k]$; one 512-bit AXI beat carries four samples
(\code{LANES=4}). We verified this packing two independent ways before
any hardware run: a bit-level reimplementation in plain \code{uint64\_t}
arithmetic checked against the \code{ap\_uint}-based packer across
4{,}000 randomized cases (including the samples whose 15-bit fields
straddle the 64-bit boundary, index 4 of 7), and a live packet capture
on loopback, decoded field-by-field against an independently
reimplemented linear congruential generator, with the high padding bits
of each 128-bit record checked for zero.

\subsection{Composition and placement}

The \code{v++} link composes four kernel instances and constrains their
connectivity:

\begin{center}
\begin{lstlisting}
fibre -> cmac_1 -> networklayer_1 -> krnl_san_scan_net_1 -> DDR -> host
                        ^
                   krnl_mm2s_1   (closes the unused TX-from-application
                                  path; v++ requires every declared
                                  stream port connected, even one this
                                  variant does not use for output)
\end{lstlisting}
\end{center}

We use the \emph{second} CMAC/network-layer instance pair
(\code{cmac\_1}, \code{networklayer\_1}), not the first: register-level
inspection via the PCIe BAR showed the first QSFP cage on our card is
uncabled. Discovering this by direct register read rather than assuming
index~0, and subsequently addressing every core by name via
\code{xrt::ip} rather than by an assumed base address, avoided a
multi-hour misconfiguration during bring-up.

The compute kernel is constrained to SLR3
(\code{slr=krnl\_san\_scan\_net\_1:SLR3}), the least-occupied
super-logic region after VNx's own footprint (VNx occupies roughly 16\%
of SLR1 and 20\% of SLR2; SLR3 sits under 1\% before our kernel is
added). Our kernel adds 11{,}299 LUTs (0.65\% of the device), 6{,}195
FFs, and 8 BRAMs, at \code{II=1} and a 3-cycle iteration latency
(\code{csynth}: Fmax 342.47~MHz).

\subsection{Timing closure}

At the requested 250~MHz kernel clock, place-and-route did \textbf{not}
close timing: worst negative slack was $-0.928$~ns on the
\code{clk\_kernel\_00} domain (2{,}508 of 212{,}143 endpoints failing,
all among a much larger sea of passing paths; the ten most-negative
individual paths in the routed timing report all terminate in
\code{krnl\_mm2s\_1}'s AXI-Lite control interface --- a stub kernel we
instantiate only to satisfy VNx's stream connectivity requirement, not
our compute path --- though the timing report does not itemize enough
endpoints to rule out other contributors). Vivado's automatic frequency
scaling reduced the kernel clock to 234.6~MHz, which is where the
implemented design actually closes. This is not a throughput-limiting
compromise for this application: at 512 bits per beat, 234.6~MHz
sustains 120.1~Gbit/s of \emph{kernel} bus bandwidth, above the
100~Gbit/s Ethernet line rate that ultimately bounds this system.

\section{Correctness}
\label{sec:correctness}

\subsection{Method}

We built a two-process test harness. \code{ctl\_san}, running on the
host that owns the U250 (over XRT/PCIe, not over the network), programs
the CMAC (enables RS-FEC, resets, and polls alignment --- link-up was
consistently observed within 0.5~s of reset in every run), writes the
network layer's identity (MAC/IP/gateway) and a single valid UDP
socket-table entry, arms the kernel with a host-computed lookup table
and the run parameters, and --- critically --- embeds the same
golden-model computation the kernel is supposed to perform, so that
every run self-verifies rather than relying on visual inspection of
printed histograms. \code{inject\_san}, running on a separate cluster
node reachable over the same fabric, generates a cohort with a
deterministic linear congruential generator and transmits it as UDP
datagrams, one packet per \code{sendto} system call
(Section~\ref{sec:throughput} explains why this specific choice --- one
packet per syscall, no batching --- is load-bearing).

\subsection{Results}

The kernel argument layout --- read from the linked \code{.xclbin}'s
embedded kernel descriptor via \code{xclbinutil --info}, not assumed ---
places the stream argument at index~0 and the scalar/buffer arguments at
indices 1, 2, 3, 4, 5, 6, 7 for the cost LUT, \code{q\_delta},
\code{n\_points}, \code{n\_samples}, histogram output, catastrophe
count, and FLOP total respectively; the lookup table is 8 entries of
64-bit MAC-prefix costs and the histogram is 8 bins --- both values
taken from the kernel source rather than assumed, after an initial
control-host implementation with a 32-bit, \code{n\_conf}-sized LUT and
a 256-bin histogram would have silently produced wrong numbers had it
not been caught before the first hardware run.

Across the full campaign --- an initial pair of runs at 100{,}003 and
4{,}000{,}003 samples, further runs during throughput characterization
(Section~\ref{sec:throughput}), and a final confirmation batch of ten
consecutive, independent 4{,}000{,}003-sample runs --- every run in
which the UDP payload was not intentionally overwhelmed
(Section~\ref{sec:fifo}) reproduced the golden model's histogram,
catastrophe count, and FLOP total exactly. The confirmation batch is
representative:

\begin{center}
\begin{lstlisting}
n_samples = 4,000,003 (x 10 independent runs)
catastrophes: 711,634 (all 10 runs)
FLOPs:        67,143,502,000,000 (all 10 runs)
histogram:    1,001,101 / 749,117 / 562,124 / 421,732 / 316,871 /
              237,424 / 711,634 / 0   (all 10 runs)
packet loss:  0 (network-layer frame counter delta across the 10 runs:
              7,148-7,170 against 7,143 data packets sent per run --
              consistent with a small amount of ordinary background
              broadcast/ARP traffic sharing the link, not data loss)
\end{lstlisting}
\end{center}

We regard this as sufficient evidence that the streaming interface
change (Section~\ref{sec:correctness}) introduces no semantic divergence
from the DMA-attached, previously-accepted kernel.

\section{Throughput: an instability that was not where we thought it was}
\label{sec:throughput}

\subsection{A genuine hardware limit, quickly found}
\label{sec:fifo}

\begin{sloppypar}
Single-packet-per-syscall UDP transmission (\code{sendto}, $\sim$104
Msamples/s natural host rate) worked reliably in our first several
hardware runs. To raise throughput, we batched transmission with
\code{sendmmsg}, submitting 64 datagrams per system call. Send-side
throughput rose to 146.6~Msamples/s --- but 97.7\% of the resulting
packets never reached the kernel (a network frame counter increased by
only 167 of the 7{,}143 packets sent). This is a real and unsurprising
finding: VNx's \code{networklayer} has small on-chip FIFOs and no
explicit flow-control back to the sender; a burst of 64 datagrams issued
by the OS in immediate succession exceeds what those FIFOs can absorb,
and excess frames are silently dropped. We do not consider this specific
finding surprising or noteworthy on its own --- it is the ordinary
behavior of a UDP/IP stack sized for moderate, not bursty, offered load
--- but its \emph{consequence} was: after this packet loss, the FPGA's
network path stopped accepting \textbf{any} further traffic, including
from a reverted, previously-reliable single-packet sender, and including
trivial 8-sample test cohorts. Reprogramming only the reconfigurable
region (\code{xrt-smi program}) did not clear this state. Recovery
required reloading the platform's \emph{base} bitstream via the card's
persistence service (which additionally re-runs CMAC bring-up) before
reprogramming our variant --- a several-minute round trip we needed six
times over the course of this work.
\end{sloppypar}

\subsection{The instability that was not the FIFO}

Having removed \code{sendmmsg} entirely (reverting to one packet per
\code{sendto} call, with an optional inter-call delay implemented as a
token bucket to explicitly pace below the batching-induced failure mode
above), we attempted a rate sweep to find a sustainable ceiling. The
results did not describe a clean threshold:

\begin{table}[h]
\centering
\small
\begin{tabular}{@{}p{0.46\textwidth}p{0.14\textwidth}p{0.28\textwidth}@{}}
\toprule
Attempt & Rate & Outcome \\
\midrule
Unpaced \code{sendto} (natural host rate) & 12.6--13.4 Gbit/s & correct in each attempt at this rate \\
Unpaced \code{sendto} (natural rate, varies with scheduling) & 17.2 Gbit/s & \textbf{failed} \\
Paced \code{sendto} & 8.0 Gbit/s & correct, 2 consecutive runs \\
Paced \code{sendto} & 12.0 Gbit/s & correct, 1 run \\
Paced \code{sendto}, same rate as the first successful pair & 8.0 Gbit/s & \textbf{failed} \\
\bottomrule
\end{tabular}
\caption{Initial rate sweep, before the ARP fix (Section~\ref{sec:arp}). The same
code, cohort, and send pattern succeeded twice and failed once at 8.0 Gbit/s.}
\end{table}

The same code, the same cohort, and the same send pattern succeeded
twice and then failed at a rate previously confirmed safe. This ruled
out a simple offered-rate threshold as the explanation and initially
read as intermittent hardware or third-party-IP misbehavior --- exactly
the diagnosis we were reluctant to accept without direct evidence,
because it is the diagnosis that most readily excuses not looking
further.

\paragraph{Root cause.}
\label{sec:arp}
\code{ip neigh show} on the sending host showed the FPGA's
neighbor-table entry in \code{STALE} state after any idle interval
between runs (editing a script, a \code{git commit}, or simply the
wall-clock gap introduced by our own test orchestration). We confirmed
this mattered operationally, not just cosmetically: five isolated
\code{ping -c1} probes, spaced one second apart, lost roughly one in
four; a single continuous run of forty pings at 0.3~s intervals lost
none. This is consistent with Linux's neighbor-discovery cache requiring
reachability reconfirmation after its reachable-time window elapses, and
with a UDP send issued while that reconfirmation is outstanding being
more likely to be delayed or dropped than one issued into an
already-warm entry --- a plausible, though not exhaustively instrumented,
mechanism, since we did not capture kernel neighbor-subsystem tracing to
confirm the exact drop point.

We also found, in the course of diagnosing this, an unrelated
orchestration bug that had been actively masking the ARP hypothesis: our
test harness used \code{set -e} and started the FPGA-side controller
process in the background before the ARP-sensitive step; when that later
step failed, the shell script exited without terminating the
still-running, still-waiting controller process, which then held the
device's exclusive context open. The \emph{next} invocation failed
immediately with an unrelated XRT error
(\code{failed to open cu context: Invalid argument}) that looked like a
fresh, distinct failure rather than a downstream symptom of the first
one. We mention this because it is a common shape of compounding failure
in iterative hardware bring-up --- a real bug's symptom is overwritten by
an orchestration bug's symptom before the real bug is understood --- and
because the fix (a \code{trap}-based cleanup that terminates an
armed-but-orphaned controller process on any subsequent failure) is
generic and reusable beyond this project.

\paragraph{Fix and validation.}
We added a three-packet ICMP warm-up (\code{ping -c3}, 0.2~s spacing)
immediately before every data burst, forcing bidirectional reachability
confirmation (the echo reply confirms the sender's neighbor entry;
producing that reply requires the FPGA's ARP responder to have the
sender resolved as well) before UDP traffic begins. We validated this
fix under conditions at least as adverse as those that triggered the
original failures: four runs immediately following the fix, the last
with the sending host's neighbor-table entry for the FPGA
\emph{forcibly deleted} beforehand (\code{ip neigh del}) to guarantee a
cold-start resolution rather than merely a stale one --- all four passed
bit-exact with zero data loss. We then ran a ten-run confirmation batch
at 4{,}000{,}003 samples and 8~Gbit/s with no manual intervention
between runs beyond the automated warm-up: \textbf{10/10 passed}, with
zero recoveries needed, reported in full in
Section~\ref{sec:correctness}.

\subsection{What we can and cannot claim about throughput}

We report 8~Gbit/s ($\approx$62.5 Msamples/s) as a \textbf{sustained,
reproducible} network-path throughput, evidenced by ten consecutive
clean runs with no failures and no recovery interventions. We do not
claim this is the network path's ceiling --- 12~Gbit/s succeeded in a
single run before the ARP-staleness confound was identified and
controlled for, and the theoretical bounds are considerably higher
(100~Gbit/s Ethernet line rate; 120.1~Gbit/s of kernel bus bandwidth at
the implemented 234.6~MHz clock) --- only that 8~Gbit/s is the rate at
which we have direct, repeated, controlled evidence of reliability.
Determining the true ceiling would require re-running the rate sweep of
Section~\ref{sec:throughput} with the ARP fix in place throughout, which
we have not yet done; we flag this as the immediate next step rather
than extrapolate it. For comparison, the previously-accepted
DMA-attached variant of the same kernel sustains 511~Msamples/s
($\approx$65.4 Gbit/s of raw data) --- the network-attached path, at its
currently-confirmed rate, is roughly 8$\times$ slower than DMA. Whether
that gap closes with a faster host-side sender (the confirmed ceiling
was bound by single-threaded \code{sendto} syscall throughput, not by
the fabric or the kernel) or persists as a structural cost of network
attachment is open.

\section{Discussion}

The central methodological point of this work is not the specific
throughput number, which we expect to improve, but the diagnostic path
that got us from ``the third-party network core is unreliable'' to
``the host's ARP cache was stale.'' Every piece of evidence available at
the network layer --- packet loss that did not correlate cleanly with
offered rate, a device that stopped responding to trivial cohorts after
a failure, recovery only via full bitstream reload --- was individually
consistent with a hardware or third-party-IP defect, and none of it, on
its own, pointed at the host's neighbor-discovery cache. What did point
there was a narrow, falsifiable observation (isolated pings losing at a
markedly higher rate than a continuous ping stream) obtained by testing
the simplest possible mechanism first (ICMP, no FPGA logic involved)
once the failure pattern (works right after another run, fails after a
gap) suggested \emph{time since last traffic}, not \emph{offered rate},
as the operative variable. We think this generalizes: when a
network-attached accelerator's failure signature correlates with
wall-clock gaps between tests rather than with load, the host IP
stack's own caches are worth checking before the accelerator's logic.

\section{Threats to validity and future work}

\begin{itemize}
  \item The root-cause mechanism for ARP staleness causing UDP loss is
    supported by a strong behavioral correlation (isolated-ping loss vs.\
    continuous-ping non-loss; failure history correlating with
    wall-clock gaps) and by the fix eliminating the failure across five
    subsequent test batches, but we did not directly instrument the
    kernel's neighbor subsystem or capture the dropped frames at the
    point of loss, so an additional or interacting factor cannot be
    fully excluded.
  \item The confirmed-safe 8~Gbit/s figure reflects ten runs at that
    specific rate; it is not a swept characterization with the ARP fix
    applied throughout, and the true throughput ceiling of the network
    path --- and whether it is bound by the host sender, the fabric, or
    the third-party network core --- remains open.
  \item All hardware validation in this work used a single Alveo U250
    on a shared cluster fabric; we have not tested a second card or a
    different network topology, and cannot yet distinguish
    card-and-fabric-specific behavior from properties of the VNx IP
    core in general.
  \item Latency (as opposed to throughput) of the network-attached path
    relative to the DMA-attached baseline is not measured in this work
    and is the more direct motivation for network attachment in a
    disaggregated pipeline; we defer it to a follow-up measurement now
    that a reliable transmission method exists.
\end{itemize}

\section*{Reproducibility}

All source --- the ported kernel, HLS testbench, the two-process host
controller and packet injector with an embedded golden model, two
independent wire-format verifiers, and orchestration scripts including
the ten-run confirmation batch --- is available in \code{fpga/san-net/}
in this repository, alongside a complete, unedited build and debugging
log in \code{fpga/san-net/README.md} covering every failure and fix
described in Section~\ref{sec:throughput}, timestamped as they occurred.

\end{document}
